\section{Data Understanding}\label{data-understanding}

This chapter covers the CRISP-DM data understanding phase for the frozen analytics base table (ABT). The goal here is simple: describe what the data looks like, where it is strong, and where it still has issues. This chapter does not discuss preprocessing or modeling steps. Those are covered in Chapter 5 and later chapters.

\subsection{\texorpdfstring{\textbf{Initial Data Collection}}{Initial Data Collection}}\label{initial-data-collection}

The final ABT had 2{,}047 rows and 50 columns. It combined three segments in one file: open-market listings, bank real and other properties acquired (ROPA), and floor-price records. This setup made it easier to compare sources before applying the training filters used later in the workflow.

Table~\ref{tab:abt-source-composition} shows the source mix. Lamudi had 1{,}619 rows (79.09\%), so it dominated the file. Bank sources had 320 rows (15.63\%). Most of these came from Bank of Commerce. BDO and Pag-IBIG had 108 floor-price rows (5.28\%).

\begin{table}[htbp]
\caption{Final ABT Composition by Source and Market Segment}
\label{tab:abt-source-composition}
\begin{center}
\scriptsize
\setlength{\tabcolsep}{4pt}
\begin{tabular}{p{0.22\textwidth} p{0.15\textwidth} p{0.09\textwidth} p{0.09\textwidth} p{0.25\textwidth}}
\toprule
Source & Market Segment & Rows & Share & Role in ABT \\
\midrule
Lamudi & open\_market & 1{,}619 & 79.09\% & Arm's-length asking-price listings used for the deployed market-value surface \\
Bank of Commerce & bank\_ropa & 234 & 11.28\% & Distressed residential inventory \\
BPI & bank\_ropa & 65 & 3.13\% & Distressed residential inventory \\
Metrobank & bank\_ropa & 15 & 0.72\% & Distressed residential inventory \\
China Bank Savings & bank\_ropa & 4 & 0.19\% & Distressed residential inventory \\
Landbank & bank\_ropa & 2 & 0.10\% & Distressed residential inventory \\
PagIBIG & floor\_price & 96 & 4.63\% & Administrative floor-price benchmark \\
BDO & floor\_price & 12 & 0.58\% & Administrative floor-price benchmark \\
\midrule
Total & all segments & 2{,}047 & 100.00\% & Frozen ABT before Chapter 5 filtering \\
\bottomrule
\end{tabular}
\end{center}
\end{table}

The ABT covered six LGUs: Cebu City, Mandaue City, Lapu-Lapu City, Talisay City, Minglanilla, and Consolacion. Naga City stayed as a CBD distance anchor, but it was not included as a training LGU. Based on the logged decision, the four Naga listings from Phase C were dropped before merge because they were too few for stable learning.


Cebu City and Lapu-Lapu City together accounted for 66.73\% of ABT rows. All 643 Lapu-Lapu rows carried \emph{is\_mactan\_island} = 1. The full breakdown by LGU and market segment is shown in Appendix B.

\begin{center}\rule{0.5\linewidth}{0.5pt}\end{center}

\subsection{\texorpdfstring{\textbf{Data Description (Schema)}}{Data Description (Schema)}}\label{data-description-schema}

The ABT columns can be grouped into eight blocks. The table below gives the main groups and their purpose.

\begin{table}[htbp]
\caption{Analytics Base Table Schema by Feature Group}
\label{tab:abt-schema}
\begin{center}
\scriptsize
\setlength{\tabcolsep}{4pt}
\begin{tabular}{p{0.18\textwidth} p{0.29\textwidth} p{0.07\textwidth} p{0.26\textwidth}}
\toprule
Feature Group & Representative Columns & Count & Role in Chapter 4 description \\
\midrule
Identifiers and provenance & property\_id, source, price\_type, geocode\_source, market\_segment & 7 & Row identity and source provenance \\
Structural and flags & property\_type, area\_sqm, floor\_area\_sqm, lot\_area\_sqm, bedrooms, bathrooms, flag columns & 10 & Physical attributes and structural sparsity \\
Location and scope & city, latitude, longitude, barangay\_geocoded, is\_mactan\_island & 5 & Study-area placement and island tagging \\
CBD-distance features & Eight retained dist variables for CBP, Mandaue, Mactan, SRP, Tabunok, Consolacion, Naga, and Airport & 8 & Network access to polycentric anchors \\
Administrative zonal fields & bir\_zonal\_rr\_median, bir\_zonal\_cr\_median, bir\_zonal\_rc\_median, bir\_zonal\_rr\_log & 4 & BIR benchmark linkage \\
Target and derived price fields & price\_php, price\_per\_sqm, log\_price, valuation\_gap, price\_outlier\_flag & 5 & Target, transform, and diagnostics \\
Spatial context & spatial\_lag\_price & 1 & Neighbor-price context \\
MCRAI accessibility fields & Nine category scores plus mcrai\_composite & 10 & Accessibility coverage summary \\
\midrule
Total & Full ABT schema & 50 & Frozen chapter reference file \\
\bottomrule
\end{tabular}
\end{center}
\end{table}

Three columns are key for this thesis flow. First, price\_per\_sqm is the target variable. Second, valuation\_gap is a diagnostic field based on price\_per\_sqm and BIR residential benchmarks. Third, the MCRAI block keeps both category-level scores and the mcrai\_composite score.

This schema also shows why the ABT was not yet a final model matrix at this stage. Some fields were complete and stable, like CBD distances and most geospatial features. Other fields depended heavily on source behavior, especially structural housing fields.

\begin{center}\rule{0.5\linewidth}{0.5pt}\end{center}

\subsection{\texorpdfstring{\textbf{Initial Data Exploration}}{Initial Data Exploration}}\label{initial-data-exploration}

Initial exploration focused on target distribution, property-type mix, geographic spread, and accessibility coverage. Out of 2{,}047 rows, 1{,}940 had non-null price\_per\_sqm. The 107 missing unit-price rows all came from Lamudi records with incomplete input fields needed for unit-price calculation.

\begin{table}[htbp]
\caption{Distribution of price\_per\_sqm by Market Segment}
\label{tab:price-by-segment}
\begin{center}
\scriptsize
\setlength{\tabcolsep}{4pt}
\begin{tabular}{p{0.16\textwidth} p{0.10\textwidth} p{0.14\textwidth} p{0.14\textwidth} p{0.16\textwidth} p{0.12\textwidth}}
\toprule
Market Segment & Non-null Rows & Q1 & Median & Mean & Q3 \\
\midrule
open\_market & 1{,}512 & PHP~65{,}323.96 & PHP~111{,}303.56 & PHP~140{,}135.47 & PHP~178{,}751.00 \\
bank\_ropa & 320 & PHP~16{,}600.00 & PHP~21{,}000.00 & PHP~30{,}000.56 & PHP~32{,}960.20 \\
floor\_price & 108 & PHP~58{,}506.38 & PHP~91{,}017.27 & PHP~90{,}188.61 & PHP~119{,}481.21 \\
\bottomrule
\end{tabular}
\end{center}
\end{table}

Table~\ref{tab:price-by-segment} shows a clear spread across segments. Open-market values were much higher than bank ROPA values on both median and mean. The open-market distribution was also right-skewed, with a long upper tail.

\begin{figure}[htbp]
\centering
\includegraphics[width=0.85\textwidth]{../../EDA/price_by_segment.png}
\caption{Distribution of price\_per\_sqm by Market Segment. Open-market values were substantially higher than bank ROPA values, with a wider upper tail reflecting the spread in residential asking prices across Metro Cebu.}
\label{fig:price-by-segment}
\end{figure}

\begin{table}[htbp]
\caption{Property-Type Composition and Median price\_per\_sqm}
\label{tab:property-type-composition}
\begin{center}
\begin{tabular}{p{0.20\textwidth} p{0.11\textwidth} p{0.11\textwidth} p{0.20\textwidth}}
\toprule
Property Type & Rows & Share & Median price\_per\_sqm \\
\midrule
Condominium & 862 & 42.11\% & PHP~166{,}666.67 \\
Vacant Lot & 461 & 22.52\% & PHP~18{,}041.24 \\
House and Lot & 412 & 20.13\% & PHP~73{,}252.03 \\
Single Detached & 221 & 10.80\% & PHP~87{,}500.00 \\
Townhouse & 89 & 4.35\% & PHP~51{,}546.54 \\
Apartment & 2 & 0.10\% & PHP~160{,}363.37 \\
\bottomrule
\end{tabular}
\end{center}
\end{table}

Table~\ref{tab:property-type-composition} shows that condominiums were the largest class in the ABT. Vacant lots had much lower median unit prices than vertical housing types. This mix helped explain the wide spread in the target field.

\begin{figure}[htbp]
\centering
\includegraphics[width=0.85\textwidth]{../../EDA/price_by_property_type.png}
\caption{Median price\_per\_sqm by Property Type. Condominiums commanded the highest unit prices; vacant lots had the lowest, reflecting their land-only composition without built area.}
\label{fig:price-by-property-type}
\end{figure}

Lapu-Lapu records also stood out in the geographic review. The city had 558 open-market rows, and all 643 rows from that LGU had is\_mactan\_island = 1. This confirms that a large part of the ABT sits in the island market context.

Geographic spread in the open-market subset was also uneven. Cebu City had the highest median unit price at PHP~142{,}886.60/sqm, followed by Lapu-Lapu City at PHP~125{,}799.56/sqm and Mandaue City at PHP~114{,}500.00/sqm. Minglanilla, Talisay City, and Consolacion formed a lower-price band, with medians between PHP~54{,}881.34 and PHP~70{,}520.26/sqm. This indicates that the premium open-market signal in the ABT remained concentrated in the core and island urban submarkets rather than the peripheral suburban LGUs. The CBD-distance matrix also showed that the southern corridor distances moved together very closely, especially the SRP--Tabunok pair ($r = 0.959$) and the Tabunok--Naga pair ($r = 0.973$). This was an early warning that the linear benchmark model would later require a collinearity check even though the nodes remained substantively meaningful.

\begin{figure}[htbp]
\centering
\includegraphics[width=0.85\textwidth]{../../EDA/price_by_city_open_market.png}
\caption{Median price\_per\_sqm by LGU for the Open-Market Subset. Cebu City and Lapu-Lapu City showed the highest median unit prices; Consolacion, Talisay, and Minglanilla formed a lower-price band.}
\label{fig:price-by-city}
\end{figure}

\begin{table}[htbp]
\caption{Overall Zero Rates for Metro Cebu Residential Accessibility Index Fields}
\label{tab:mcrai-zero-rates}
\begin{center}
\begin{tabular}{p{0.23\textwidth} p{0.13\textwidth} p{0.44\textwidth}}
\toprule
MCRAI Field & Zero Rate & Exploratory Reading \\
\midrule
mcrai\_education & 39.86\% & Education access remained sparse for peripheral and island locations \\
mcrai\_finance & 36.34\% & Finance nodes were still clustered around denser urban centers \\
mcrai\_health & 32.10\% & Hospital and clinic access was uneven across the six-LGU scope \\
mcrai\_security & 30.22\% & Security infrastructure remained absent in many residential catchments \\
mcrai\_retail\_density & 21.59\% & Retail and food density improved after query expansion but was not universal \\
mcrai\_recreation & 6.80\% & Recreation coverage was broad but still thinner in selected peripheral locations \\
mcrai\_transport & 5.59\% & Corridor-based transport access was widespread, not complete \\
mcrai\_grocery & 2.60\% & Grocery access was near-universal in the final file \\
mcrai\_tourism & 0.58\% & Tourism and hospitality nodes were almost always observed within the broad catchment \\
mcrai\_composite & 0.24\% & Positive-weight composite coverage was effectively complete \\
\bottomrule
\end{tabular}
\end{center}
\end{table}

Table~\ref{tab:mcrai-zero-rates} shows that MCRAI coverage was generally strong, but not uniform across categories. Grocery, tourism, and the composite score had very low zero rates. Education, finance, health, and security had higher zero rates.

A rank-based exploratory check on the open-market subset showed that location and accessibility fields aligned positively with unit price, while several size-related structural fields moved in the opposite direction. \texttt{bir\_zonal\_rr\_median} ($\rho = 0.343$) and \texttt{mcrai\_composite} ($\rho = 0.273$) were positively associated with \texttt{price\_per\_sqm}. By contrast, \texttt{floor\_area\_sqm} ($\rho = -0.461$), \texttt{bedrooms} ($\rho = -0.450$), and \texttt{bathrooms} ($\rho = -0.321$) were negatively associated with \texttt{price\_per\_sqm}. This does not mean the target was inverted. It reflects the arithmetic of a unit-price target and the composition of the sample: compact condominiums in Cebu City and Lapu-Lapu City tend to command higher PHP/sqm, while larger detached homes and land-oriented listings more often sit in lower-density submarkets where total asking price can still be high but price per square meter is lower. Consistent with this reading, the same structural variables remained positively associated with total asking price in the open-market subset.

The updated retail-density behavior after the targeted POI re-fetch is visible in the frozen ABT. Minglanilla had 6.74\% overall zeros in mcrai\_retail\_density, and 5.63\% in the open-market subset. In the same open-market subset, Talisay City had 18.82\% and Lapu-Lapu City had 24.33\%. For transport, open-market zero rates remained higher in Minglanilla (29.58\%) and Lapu-Lapu City (14.85\%).

For BIR linkage, the frozen ABT had full coverage for bir\_zonal\_rr\_median across all 2{,}047 rows. The incomplete part was valuation\_gap, which was missing in the same 107 rows where price\_per\_sqm was missing.

\begin{center}\rule{0.5\linewidth}{0.5pt}\end{center}

\subsection{\texorpdfstring{\textbf{Data Quality Assessment}}{Data Quality Assessment}}\label{data-quality-assessment}

Data quality checks looked at missingness, integrity flags, cross-segment differences, and external-data contamination issues that were fixed before freeze.

\begin{table}[htbp]
\caption{Major Missingness Patterns in the Final ABT}
\label{tab:missingness-patterns}
\begin{center}
\scriptsize
\setlength{\tabcolsep}{4pt}
\begin{tabular}{p{0.18\textwidth} p{0.09\textwidth} p{0.09\textwidth} p{0.20\textwidth} p{0.22\textwidth}}
\toprule
Field & Nulls & Share & Concentration & Reading for Chapter 5 \\
\midrule
lot\_area\_sqm & 1{,}758 & 84.72\% & 100.00\% of open\_market rows; 82.41\% of floor\_price rows & Structural source absence, not random missingness \\
barangay\_geocoded & 834 & 40.19\% & Weakest barangay recovery in Consolacion and Mandaue & Location granularity remained uneven \\
bedrooms & 334 & 16.10\% & 20.28\% of open\_market rows & Mostly Lamudi field gaps \\
bathrooms & 237 & 11.42\% & 14.39\% of open\_market rows & Mostly Lamudi field gaps \\
price\_per\_sqm & 107 & 5.23\% & All 107 nulls came from Lamudi rows & Target derivation incomplete for a small open-market subset \\
valuation\_gap & 107 & 5.23\% & Same rows as price\_per\_sqm & Diagnostic field depends on target availability \\
floor\_area\_sqm & 78 & 3.76\% & 4.74\% of open\_market rows & Residual Lamudi incompleteness \\
area\_sqm and log\_price & 78 each & 3.76\% & Same rows as floor\_area\_sqm nulls & Derived-field propagation of source gaps \\
spatial\_lag\_price & 23 & 1.11\% & Thinly spread across all segments & Neighbor-based context unavailable for a small edge subset \\
\bottomrule
\end{tabular}
\end{center}
\end{table}

The largest missingness issue was structural. lot\_area\_sqm was fully missing in open-market rows and heavily missing in floor-price rows. This reflects source limitations, not random noise.

\begin{figure}[htbp]
\centering
\includegraphics[width=0.85\textwidth]{../../EDA/missingness_top15.png}
\caption{Top-15 Fields by Missingness Rate in the Frozen ABT. lot\_area\_sqm dominated because it was structurally absent in virtually all open-market rows rather than randomly missing.}
\label{fig:missingness}
\end{figure}

The remaining structural gaps were mostly in Lamudi rows. In the frozen ABT, floor\_area\_sqm had 78 nulls total, and bedrooms and bathrooms nulls were concentrated in the open-market segment.

The ABT also kept soft outlier flags separate from hard integrity drops. There were 19 rows with price\_outlier\_flag = True (14 open-market, four bank ROPA, one floor-price). These flags did not automatically remove rows. A smaller integrity set was identified for downstream removal: row 431 (property\_id 468), row 658 (property\_id 714), and row 705 (property\_id 769). Their unit prices were PHP~3.68, PHP~4.13, and PHP~4.60, which are not plausible for residential valuation in this context.

Cross-segment behavior was also clear. Bank ROPA unit prices were systematically lower than open-market values. Floor-price records represented a different pricing basis. This supports the later segmentation step tied to valuation standards, where open\_market is used as the modeling scope for Market Value estimation \parencite{ivs2025}.

Another issue was fixed before freeze in the POI pipeline. Earlier runs had a bounding-box bug that allowed out-of-scope points (for example from Bohol, Batangas, and Tagaytay) to pass the LGU filter. The corrected freeze removed this contamination before final MCRAI recomputation.

Overall, the ABT at this stage was rich and usable for analysis, but still not model-ready in raw form. It had strong geospatial and benchmark coverage, but it also had source-driven structural gaps, a few implausible rows, and clear segment differences in price behavior. Chapter 5 explains how these findings were handled in data preparation.
